\begin{lemma}[structure-sheaf cohomology of the cusp fibre]
\label{lem:cusp-structure-sheaf-audit}
\Verified\ The reduced normal-crossings fibre \(W\) satisfies
\[
 h^q(W,\mathcal O_W)=(1,2,1)\qquad(q=0,1,2).
\tag{9.7}
\]
The natural map
\[
 H^1(W;\mathbb C)\longrightarrow H^1(W,\mathcal O_W)
\tag{9.8}
\]
is an isomorphism.  In degree two the corresponding map is surjective and
is injective on the kernel of pullback to the normalization
\(\nu:\mathrm{dP}_6\to W\).  Finally,
\[
 \bigwedge^2H^1(W,\mathcal O_W)
 \longrightarrow H^2(W,\mathcal O_W)
\tag{9.9}
\]
is an isomorphism.
\end{lemma}

\begin{proof}
Let \(D_1,D_2,D_3\simeq\mathbb P^1\) be the double curves and \(P,Q\) the
two triple points, with the branch order fixed in the proof of
\cref{thm:p0-integral-specialization}.  The normalization complex is
\[
0\longrightarrow\mathcal O_W\longrightarrow
 \nu_*\mathcal O_{\mathrm{dP}_6}\xrightarrow{\delta_0}
 \bigoplus_{k=1}^3\mathcal O_{D_k}\xrightarrow{\delta_1}
 \mathbb C_P\oplus\mathbb C_Q\longrightarrow0,
\tag{9.10}
\]
where \(\delta_0\) is the difference of the two branch restrictions and
\[
 \delta_1(g_1,g_2,g_3)
 =(g_1(P)-g_2(P)+g_3(P),
   g_1(Q)-g_2(Q)+g_3(Q)).
\]
For completeness, exactness is stalkwise.  At a generic double point it is
the fibre-product sequence
\[
0\to\mathbb C\{x,y,z\}/(xy)
 \to\mathbb C\{y,z\}\oplus\mathbb C\{x,z\}
 \to\mathbb C\{z\}\to0.
\]
At a triple point it is the augmented branch complex for
\(\mathbb C\{x,y,z\}/(xyz)\): three branch rings, three pairwise
intersection rings, and the common residue field, with alternating
differences.  The augmented two-simplex complex is exact; successively
lifting compatible power series proves exactness of the analytic rings.
Away from the singular locus the first arrow is an isomorphism.

The same formulas give an exact complex with \(\mathbb C\) in place of
\(\mathcal O\).  Since \(\mathrm{dP}_6\) is rational and each \(D_k\) is a
rational curve, the positive structure-sheaf cohomology of every term on
the right of (9.10) vanishes.  Hypercohomology therefore reduces to
\[
 \mathbb C\xrightarrow{0}\mathbb C^3
 \xrightarrow{d_1}\mathbb C^2,
 \qquad
 d_1(a_1,a_2,a_3)=(s,s),\quad s=a_1-a_2+a_3.
\tag{9.11}
\]
Its cohomology has dimensions \(1,2,1\), proving (9.7).

For the constant-sheaf complex, the row of degree zero is again (9.11),
the row of degree one vanishes, and the remaining degree-two row begins
with \(H^2(\mathrm{dP}_6;\mathbb C)\).  The comparison induced by
\(\mathbb C_W\hookrightarrow\mathcal O_W\) is the identity on (9.11) and
zero on the positive rows of the target.  It follows that (9.8) is an
isomorphism, while in degree two it maps the dual-complex summand
\(\ker\nu^*\) isomorphically onto \(H^2(W,\mathcal O_W)\).

Choose the basis \(\gamma_W,u_W\in H^1(W;\mathbb Z)\) whose intrinsic
specializations are \(\gamma,u\), respectively; it exists and is unique by
\cref{thm:p0-integral-specialization}.  Their cup product is nonzero,
because its specialization is
\(\gamma\wedge u\ne0\in H^2(F;\mathbb Z)\).  Its pullback to
\(\mathrm{dP}_6\) is zero since
\(H^1(\mathrm{dP}_6;\mathbb C)=0\).  The degree-two comparison therefore
sends it to a nonzero element of \(H^2(W,\mathcal O_W)\).  The images of
\(\gamma_W,u_W\) form a basis by (9.8), and both sides of (9.9) are
one-dimensional.  Hence (9.9) is an isomorphism.
\end{proof}
